\documentclass[10pt,a4paper]{amsart}
\usepackage[margin=19mm]{geometry}
\usepackage{amsmath,amssymb,mathtools}
\usepackage[scr=rsfs,cal=euler]{mathalfa}
\usepackage{xfrac}
\usepackage[unicode,hidelinks,bookmarks=false]{hyperref}
\newcommand{\ie}{i.e.\ }
\newcommand{\cf}{cf.\ }
\newcommand{\resp}{respectively\ }
\newcommand{\Subsection}{\subsection}
\providecommand{\nobreakdash}{\nobreak\hbox{-}}
\setlength{\emergencystretch}{2em}
\multlinegap0pt
\usepackage{bm}
\usepackage{bbm}
\usepackage{dsfont}
\usepackage{xspace}
\usepackage{cancel}
\usepackage{mathtools}
\usepackage{amsthm}
\makeatletter
\@ifundefined{prp}{\newtheorem{prp}{Prp}}{}
\@ifundefined{lemp}{\newtheorem{lemp}[prp]{Lemma}}{}
\makeatother
\newcommand{\C}{\mathbb{C}}
\newcommand{\RI}{{\mathrm {I}}}
\title[Archimedean Modular Symbols for Rankin-Selberg Convolutions ]{Archimedean Modular Symbols for Rankin-Selberg Convolutions of $\mathrm{GL}\_n(\mathbb{C})\times\mathrm{GL}\_n(\mathbb{C})$ and $\mathrm{GL}\_n(\mathbb{C})\times\mathrm{GL}\_{n+1}(\mathbb{C})$}
\author{Yubo Jin, Zeyu Tian}
\date{Source: arXiv:2607.19133v1 (CC BY 4.0)}
\begin{document}
\maketitle

\noindent\emph{Source and licence.} Archimedean Modular Symbols for Rankin-Selberg Convolutions of $\mathrm{GL}_n(\mathbb{C})\times\mathrm{GL}_n(\mathbb{C})$ and $\mathrm{GL}_n(\mathbb{C})\times\mathrm{GL}_{n+1}(\mathbb{C})$. Version arXiv:2607.19133v1 (CC BY 4.0), distributed under the Creative Commons Attribution 4.0 International licence, as recorded in the arXiv licence metadata for this exact version and shown on the abstract page https://arxiv.org/abs/2607.19133v1. Full rights notice: licenses/arxiv-2607.19133-CC-BY-4.0.txt.\par\smallskip

\noindent\textit{Editorial scope.} Counted unit: the Lemp whose source label is \texttt{lemD} in the article (source0.tex lines 1220--1225, source label \texttt{lemD}), delivered together with the article's own complete original proof of it (source0.tex lines 1227--1278, one \texttt{proof} environment of 52 lines). No mathematics is written, repaired or re-derived: statement and proof are the article's own text line for line. Required context retained uncounted: lemB (lines 1140--1145); lemC (lines 1198--1203). Every definition, display and label of those lines is retained unchanged. Omitted from this package and not redistributed: 1--1139, 1146--1197, 1204--1219, 1226--1226, 1279--1502. Nothing inside a retained excerpt is omitted or altered: every retained body is delivered exactly as the article's lines are written, so no omission receipt of any kind accompanies this package. Citations: the retained excerpts carry no citation command at all, so no bibliography entry of the article is transcribed and no reference is redistributed. Reference adaptation: none was needed and none was made; every reference inside the delivered unit resolves inside it, so no unresolved reference or citation is produced. Printed slips kept literal: no printed word, symbol, hypothesis or number of the counted statement or of its proof was altered. Layout and class: the article's own class is \texttt{amsart} with packages including amsmath, latexsym, amsfonts, amssymb, amsxtra, bbm, color, mathrsfs, amscd, mathdots, braket, hyperref, pdfsync, xypic; the wrapper is the lane's pinned \texttt{amsart} editorial wrapper at 10pt on A4, which restates the theorem-like environments the retained text needs and adds the article's own macro definitions that the retained text needs (\texttt{\textbackslash C}, \texttt{\textbackslash RI}), with extra wrapper packages (bm, bbm, dsfont, xspace, cancel, mathtools). The delivered numbering is the unit's own. Assets: no retained span contains a figure environment, a graphics-inclusion command, a PDF-page inclusion, a table or a TikZ picture; no image file is copied into this package, the package declares no asset dependency and no image file is redistributed with it. Licence: version arXiv:2607.19133v1 of this article is distributed under the Creative Commons Attribution 4.0 International licence, as recorded in the arXiv licence metadata for this exact version and shown on the abstract page https://arxiv.org/abs/2607.19133v1. A full rights notice is delivered at licenses/arxiv-2607.19133-CC-BY-4.0.txt. Characters: every retained excerpt is pure ASCII, so the delivered engine, XeLaTeX with its default Latin Modern text font, reports no missing character.
	\begin{lemp}\label{lemB}
		We have that
		\[
		\mathfrak{o}_{n,n}=\omega_{H^{\vee}_n}\wedge\omega_{H_n}\wedge\left(e_{1}^{\vee}\wedge e_{2}^{\vee}\cdots\wedge e_{n-1}^{\vee}\right).
		\]
	\end{lemp}

	\begin{lemp}\label{lemC}
		We have that
		\[
		\xi_{P_{\circ}^{\vee}}=\frac{1}{n!(n-1)!\cdots 2!}\cdot(\varepsilon_{n-1,n})^1(\varepsilon_{n-2,n-1})^2\cdots(\varepsilon_{1,2})^{n-1}\xi_{H'^{\vee}_n}.
		\]
	\end{lemp}

	\begin{lemp}\label{lemD}
		We have that
		\[
		\omega_{H_n}\wedge\omega_{H_n^{\vee}}\wedge\omega'_{P_{\circ}}=\frac{(-1)^{\frac{n(n-1)}{2}}}{n!}\cdot\mathfrak{o}_{n,n}.
		\]
	\end{lemp}

	\begin{proof}
		We write
		\[
		\omega'_{P_{\circ}^{\vee}}=\sum_{\substack{
				1\leq i_1\leq i_2\leq \cdots\leq i_{n-1}\leq n-1\\
				1\leq j_1\leq j_2\leq\cdots\leq j_{n-1}\leq n
		}} C_{i_1,i_2,\dots,i_{n-1}}^{j_1,j_2,\dots,j_{n-1}}\cdot e^{\vee}_{i_1,j_1}\wedge e^{\vee}_{i_2,j_2}\wedge\cdots\wedge e^{\vee}_{i_{n-1},j_{n-1}},
		\]
		where $C_{i_1,i_2,\dots,i_{n-1}}^{j_1,j_2,\dots,j_{n-1}}\in\C$ and $e_{i,i}^{\vee}:=e_i^{\vee}$ for $1\leq i\leq n-1$. In view of Lemma \ref{lemB}, we have that
		\[
		\omega_{H_n}\wedge\omega_{H_n^{\vee}}\wedge\omega'_{P_{\circ}}=C_{1,2,\dots,n-1}^{1,2,\dots,n-1}\cdot\mathfrak{o}_{n,n}.
		\]
		
		Denote by $\varepsilon_{i,j}^{\mathfrak{k}}\in\mathfrak{k}_n$ the corresponding element of $\varepsilon_{i,j}$ ($1\leq i,j\leq n$) under the canonical identification $\mathfrak{k}_n=\mathfrak{gl}_n(\C)$. Then Lemma \ref{lemC} implies that
		\[
		\omega'_{P_{\circ}^{\vee}}=\frac{1}{n!(n-1)!\cdots 2!}\cdot(\varepsilon^{\mathfrak{k}}_{n-1,n})^1(\varepsilon^{\mathfrak{k}}_{n-2,n-1})^2\cdots(\varepsilon^{\mathfrak{k}}_{1,2})^{n-1}\omega'_{H'^{\vee}_n}.
		\]
		Recall that $\RI_n'(\Xi_n^{\vee})$ has the highest weight vector
		\[
		\omega'_{\check{H}_n'}=\omega_n'=e_{n,1}^{\vee}\wedge e_{n,2}^{\vee}\wedge\cdots\wedge e_{n,n-1}^{\vee}
		\]
		and it is easy to obtain the lowest weight vector
		\[
		\omega'_{H'^{\vee}_n}=(-1)^{n-1}e_{1,2}^{\vee}\wedge e_{1,3}^{\vee}\wedge\cdots\wedge e_{1,n}^{\vee}.
		\]
		We compute that
		\[
		\begin{aligned}
			&\,(\varepsilon^{\mathfrak{k}}_{n-1,n})^1(\varepsilon^{\mathfrak{k}}_{n-2,n-1})^2\cdots(\varepsilon^{\mathfrak{k}}_{1,2})^{n-1}\omega'_{H'^{\vee}_n}\\
			=&\,(\varepsilon^{\mathfrak{k}}_{n-1,n})^1(\varepsilon^{\mathfrak{k}}_{n-2,n-1})^2\cdots(\varepsilon^{\mathfrak{k}}_{1,2})^{n-1}\left(e_{1,2}^{\vee}\wedge e_{1,3}^{\vee}\wedge\cdots\wedge e_{1,n}^{\vee}\right)\\
			=&\,(-1)^{\frac{(n-1)(n-2)}{2}}\cdot(n-1)!(n-2)!\cdots2!\cdot\left(e_{1}^{\vee}\wedge e_{2}^{\vee}\wedge\cdots\wedge e_{n-1}^{\vee}\right)\\
			&+\text{(other terms)},
		\end{aligned}
		\]
		where `other terms' indicate the sum of elements of the form
		\[
		C\cdot e^{\vee}_{i_1,j_1}\wedge e^{\vee}_{i_2,j_2}\wedge\cdots\wedge e^{\vee}_{i_{n-1},j_{n-1}}\qquad (C\in\C^{\times})
		\]
		with
		\[
		1\leq i_1\leq i_2\leq \cdots\leq i_{n-1}\leq n-1,\qquad 1\leq j_1\leq j_2\leq\cdots\leq j_{n-1}\leq n,
		\]
		and
		\[
		\left(i_1,i_2,\dots,i_{n-1}\right)\neq\left(1,2,\dots,n-1\right),\quad\left(j_1,j_2,\dots,j_{n-1}\right)\neq\left(1,2,\dots,n-1\right).
		\]
		Hence
		\[
		C_{1,2,\dots,n-1}^{1,2,\dots,n-1}=\frac{(-1)^{\frac{n(n-1)}{2}}}{n!},
		\]
		which completes the proof of the lemma.
	\end{proof}

\end{document}
